\documentclass[11pt,a4paper]{article}
\usepackage[margin=2.3cm]{geometry}
\usepackage{amsmath,amssymb,amsthm}
\usepackage{empheq}
\usepackage{booktabs}
\usepackage{graphicx}
\usepackage{xcolor}
\usepackage{caption}
\usepackage{listings}
\usepackage{array}
\usepackage{multirow}
\usepackage[hidelinks]{hyperref}
\usepackage{microtype}
\definecolor{navy}{HTML}{1E2761}
\definecolor{dkred}{HTML}{993C1D}
\definecolor{dkgreen}{HTML}{0F6E56}
\definecolor{codebg}{HTML}{F6F7F9}
\definecolor{boxbg}{HTML}{F4F6FA}
\lstset{
  basicstyle=\ttfamily\footnotesize,
  backgroundcolor=\color{codebg},
  frame=single, rulecolor=\color{gray!40},
  breaklines=true, showstringspaces=false,
  keywordstyle=\color{navy}\bfseries,
  commentstyle=\color{dkgreen}\itshape,
  stringstyle=\color{dkred},
  language=Python, numbers=left, numberstyle=\tiny\color{gray},
  xleftmargin=1.6em, framexleftmargin=1.6em
}
\newtheorem{proposition}{Proposition}
\newcommand{\A}{\mathcal{A}}

% ---- a labelled box stating which data a section used -------------------
\newcommand{\datatag}[2]{%
  \begin{center}\small
  \colorbox{boxbg}{\parbox{0.965\linewidth}{\textbf{Data:} #1 \\ \textbf{Code:} \texttt{#2}}}
  \end{center}}
\newcommand{\verdictgood}[1]{\textcolor{dkgreen}{\textbf{#1}}}
\newcommand{\verdictbad}[1]{\textcolor{dkred}{\textbf{#1}}}

\title{\textbf{Base-Rate and Operating-Point Confounding\\in Crime-Prediction Fairness Metrics}\\[4pt]
\large A technical report: every method tried, the mathematics, and the results obtained}
\author{Gaurav Jha}
\date{\today}

\begin{document}
\maketitle

\begin{abstract}
\noindent
Reported fairness gaps in crime prediction are largely attributable to
\emph{base rates} rather than to differences in how well a model serves each
neighbourhood. We show further that correcting for base rate is
\emph{insufficient}: a model can close a base-rate-corrected gap entirely by
shifting its decision threshold while its ranking ability degrades. We
demonstrate this five independent ways --- in a federated spatio-temporal GNN,
in a threshold sweep requiring no architectural change at all, in a
physics-informed neural network (PINN) built on the Short et al.\ (2008)
burglary equations, in a learned-threshold experiment where optimising the
fairness-relevant metric \emph{creates} a $35.56$-point disparity, and in an
audit of the two most widely used fairness libraries, whose default settings
reproduce the same failure. Across four PINN configurations, real skill (AUC)
spans $18.32$ points while the Head--Tail $F_1$ gap spans $2.06$; within the
Tail group alone, $F_1$ varies by $0.28$ while AUC varies by $7.2$. The PINN is
the only intervention of seven tested that beat its capacity-matched control
($+18.32$ AUC, $t = 9.57$, five seeds). All model results use bit-reproducible
training with capacity-matched controls.
\end{abstract}

\tableofcontents

%======================================================================
\section{Where we started, and where we are now}

\subsection{The original problem}
We reproduced FedCrime (Neurocomputing 679, 2026): federated crime prediction
across six police stations with a zero-inflated negative binomial (ZINB) loss.

\subsection{Where we are now}
The gap is largely \emph{not a property of the model}. Section~\ref{sec:theory}
gives the mathematics; Sections~\ref{sec:phases1}--\ref{sec:tools} give each
method together with the results it produced. Seven interventions were tested
against capacity-matched controls; six failed, and the one that succeeded
(the PINN) simultaneously provided the cleanest demonstration that the standard
fairness metric does not track model quality.

\subsection{A note on which data each result uses}
This project uses real data for its empirical findings and purpose-built
synthetic data for two \emph{controlled} experiments. The distinction is
important and is stated at the head of every section.

\begin{center}
\begin{tabular}{lll}
\toprule
\textbf{Section} & \textbf{Experiment} & \textbf{Data} \\
\midrule
\ref{sec:phases1} & FedCrime reproduction, federated baselines & real (LA 2018) \\
\ref{sec:diagnosis} & Base-rate diagnosis, four cities, literature audit & real \\
\ref{sec:interventions} & Reweighting, Group-DRO, adaptive thresholds & real (LA) \\
\ref{sec:graphs} & Graph architectures, depth sweep & real (LA + Chicago) \\
\ref{sec:sweep} & Threshold sweep & real (Chicago) \\
\ref{sec:pinn} & Physics-informed model & real (Chicago burglary) \\
\ref{sec:learned} & Learned thresholds, Arm A / Arm B & \textbf{synthetic} \\
\ref{sec:tools} & Fairlearn / AIF360 audit & \textbf{synthetic} \\
\bottomrule
\end{tabular}
\end{center}

\noindent
Sections~\ref{sec:learned} and \ref{sec:tools} \emph{must} use synthetic data.
Their claim is that a method reports or creates a disparity that does not exist.
Establishing that requires knowing with certainty that no disparity exists,
which on real data is unknowable --- sparse neighbourhoods may genuinely be
harder to predict. By generating the data we make equal skill a property of the
construction rather than an estimate, so any disparity the method reports is
provably an artifact. This is instrument calibration against a known standard,
not a shortcut.

%======================================================================
\section{Mathematical background}\label{sec:theory}

\subsection{The skill-free baseline of $F_1$}
Let a group have base rate $p = \Pr(y=1)$. Consider the \emph{skill-free}
predictor $\hat{y}\equiv 1$: it ignores all inputs and predicts the positive
class everywhere. Over $n$ samples,
\[
TP = np, \qquad FP = n(1-p), \qquad FN = 0 .
\]
Substituting into $F_1 = \dfrac{2TP}{2TP + FP + FN}$:
\begin{equation}
\A_{F_1}(p) \;=\; \frac{2np}{2np + n(1-p)} \;=\; \frac{2p}{1+p}.
\label{eq:af1}
\end{equation}

\begin{proposition}[Attainability is base-rate dependent]
$\A_{F_1}$ is strictly increasing in $p$ on $(0,1)$, since
$\A_{F_1}'(p) = 2/(1+p)^2 > 0$. Hence two groups with identical model
skill but different base rates yield different $F_1$, and their difference is
non-zero by construction.
\end{proposition}

\noindent Concretely, at $p=0.50$ we get $\A_{F_1} = 66.7$; at $p=0.05$,
$\A_{F_1}=9.5$. A model that never reads the data exhibits a $57$-point
``fairness gap''.

\subsection{Base-rate invariance of AUC}
For the same skill-free predictor all scores are tied, so the Mann--Whitney
statistic gives
\begin{equation}
\mathrm{AUC} \;=\; \Pr\!\big(s(x^+) > s(x^-)\big) + \tfrac12\Pr\!\big(s(x^+)=s(x^-)\big) \;=\; \tfrac12
\label{eq:auc}
\end{equation}
independently of $p$. AUC conditions on class membership, so class balance
cancels. This is why we treat AUC as the primary evidence throughout.

\subsection{The skill score, and why it is not enough}
The natural correction normalises the observed metric $M$ against its skill-free
baseline:
\begin{equation}
s \;=\; \frac{M - \A(p)}{1 - \A(p)},
\label{eq:skill}
\end{equation}
which is $0$ at the skill-free baseline and $1$ at perfect prediction.

\paragraph{The failure.} Equation~\eqref{eq:skill} removes the dependence on
$p$. It does \emph{not} remove the dependence on the decision threshold $\tau$.
Writing $F_1$ explicitly as a function of both,
\begin{equation}
F_1(p,\tau) \;=\; \frac{2\,\mathrm{TPR}(\tau)\,p}{\mathrm{TPR}(\tau)\,p + \mathrm{FPR}(\tau)(1-p) + p},
\label{eq:f1ptau}
\end{equation}
we see two free parameters. Correcting for $p$ leaves $\tau$ untouched. Since
$\mathrm{TPR}$ and $\mathrm{FPR}$ are both monotone in $\tau$, a model can
increase $F_1$ purely by lowering $\tau$ --- while
$\mathrm{AUC} = \int_0^1 \mathrm{TPR}\,d\,\mathrm{FPR}$ is unchanged, being an
integral over \emph{all} thresholds.

\medskip
\noindent\fbox{\parbox{0.96\linewidth}{%
\textbf{Central claim.} Base-rate correction of a threshold-dependent metric is
insufficient. Base rate and operating point are two distinct confounds;
removing one leaves the other. An invariant metric is necessary, not merely
preferable.}}
\medskip

\noindent We stress that Eqs.~\eqref{eq:af1}--\eqref{eq:auc} are
\emph{established} \cite{davis2006,boyd2012,chouldechova2017} and are used here
as background. The contributions are empirical.

%======================================================================
\section{Phase 1 --- Reproduction and baselines}\label{sec:phases1}
\datatag{Real. Los Angeles 2018, 113 neighbourhoods, 8 crime types, 365 days.}{src/fedcrime/}

\subsection{Method}
A three-block temporal convolutional network (TCN) with dilations $\{1,2,4\}$
and residual connections, followed by a mixing stage, then ZINB heads
$(\pi,\mu,\phi)$ and a binary classifier head. Federated training uses FedAvg
over six clients formed by round-robin assignment of regions ranked by crime
volume. Class imbalance is handled with per-category positive weights
$w_c = \mathrm{clip}(n^-_c/n^+_c, 1, 10)$:
\[
\mathcal{L} \;=\; \mathrm{BCE}_{w}(\hat{y}, y) \;+\; 0.3\,\mathcal{L}_{\mathrm{ZINB}}.
\]

\subsection{Results --- reproduction succeeded}
We recovered the published trend as the data subset became sparser:
\begin{center}
\begin{tabular}{lc}
\toprule
Subset & $F_1$ \\
\midrule
Densest & $72.16$ \\
Sparsest & $48.43$ \\
\bottomrule
\end{tabular}
\end{center}

\subsection{Results --- the disparity appears}
Partitioning regions by training-period crime volume into Head / Mid / Tail
(20/30/50 percentiles):
\begin{center}
\begin{tabular}{lccc}
\toprule
& Head & Mid & Tail \\
\midrule
$F_1$ (LA 2018) & $69.4$ & $47.0$ & $\mathbf{28.2}$ \\
\bottomrule
\end{tabular}
\end{center}
\noindent A $\mathbf{41}$-point gap. Closing it became the research question.

\subsection{Results --- federated baselines}
Nine federated aggregation methods were implemented for comparison: FedAvg,
FedProx, SCAFFOLD, MOON, FBLG, FedAvgM, FedTrimAvg, Multi-Krum and Bulyan.
A defect in the Bulyan implementation (\texttt{sset\_w} / \texttt{sel\_size})
was found and corrected during testing. No aggregation rule altered the
Head--Tail gap materially, which was the first indication that the gap was not
an optimisation artifact.

%======================================================================
\section{Phase 2 --- Diagnosis of the confound}\label{sec:diagnosis}
\datatag{Real. Four cities; six published models; FedCrime's own published tables; Chicago 311 service requests. Plus one controlled simulation, marked below.}{analysis/theory.py, baserate\_analysis.py, literature\_audit.py, reanalysis\_published.py, decision\_impact.py}

\subsection{Results --- the theory, verified}
\begin{center}
\begin{tabular}{lc}
\toprule
Check & Result \\
\midrule
$\A_{F_1}(p)=2p/(1+p)$ vs.\ empirical $F_1$ & matches, 8 base rates \\
AUC of a skill-free predictor & $0.5$ at every $p$ \\
Controlled simulation, \textbf{equal skill by construction} (synthetic) & $F_1$ gap $\mathbf{38.69}$, AUC gap $\mathbf{0.19}$ \\
\bottomrule
\end{tabular}
\end{center}
\noindent The third row is decisive: with skill held identical by construction,
$F_1$ still reports a $38.69$-point gap while AUC correctly reports $0.19$.

\subsection{Results --- four cities}
Artifact gap, i.e.\ the $F_1$ gap attributable to base rate alone:
\begin{center}
\begin{tabular}{lcccc}
\toprule
& Los Angeles & Chicago & New York & San Francisco \\
\midrule
Artifact gap & $42.8$ & $45.9$ & $19.9$ & $48.8$ \\
\bottomrule
\end{tabular}
\end{center}
\noindent Correlation with sparsity: $r = +0.835$. New York's lower value is
consistent with its less extreme spatial concentration.

\subsection{Results --- the published literature}
Correlation between reported $F_1$ and a base-rate proxy across six published
models, under six different assumptions:
\begin{center}
\begin{tabular}{lcc}
\toprule
Assumption for base rate & macro-$F_1$ & micro-$F_1$ \\
\midrule
Poisson & $0.642$ & $0.868$ \\
Negative binomial, $k=5$ & $0.651$ & $0.875$ \\
Negative binomial, $k=1$ & $0.662$ & $0.884$ \\
Raw rate (no distribution) & $0.696$ & $0.899$ \\
Log crime volume & $0.672$ & $0.890$ \\
\textbf{Rank of crime volume only} & $\mathbf{0.719}$ & $\mathbf{0.875}$ \\
\bottomrule
\end{tabular}
\end{center}
\noindent The final row assumes \emph{nothing} --- only the ordering of groups,
stated directly in the source papers --- and is \emph{higher} than the
distributional version, so the finding does not depend on our modelling choices.
Fisher combined $p = 1.32\times10^{-5}$; bootstrap $95\%$ CI over models
$[+0.694,\,+0.974]$. ST-HSL is non-significant ($p = 0.36$) and serves as an
internal negative control.

\subsection{Results --- FedCrime's own tables, and a second domain}
\begin{center}
\begin{tabular}{lc}
\toprule
Re-analysis & Correlation \\
\midrule
FedCrime's own published tables & $r = 0.998$ \\
Chicago 311 service requests (non-crime domain, same spatial units) & same pattern \\
\bottomrule
\end{tabular}
\end{center}

\subsection{Results --- does it change decisions?}
\begin{center}
\begin{tabular}{lc}
\toprule
Decision & Changes after skill correction? \\
\midrule
Which model is best, per group & $\mathbf{0}$ of $5$ groups \\
Which group needs remediation & $\mathbf{6}$ of $6$ models \\
\bottomrule
\end{tabular}
\end{center}
\noindent Mean Spearman $\rho$ between the raw-$F_1$ ranking and the
skill-corrected ranking: $\mathbf{-0.700}$. The negative half of this result
matters as much as the positive: the confound does \emph{not} invalidate model
comparisons, only remediation targeting. Overstating the scope would be an
error.

%======================================================================
\section{Phase 3 --- Fairness interventions}\label{sec:interventions}
\datatag{Real. Los Angeles.}{src/fedcrime/robust\_fair\_gnn.py}

\subsection{Reweighting}
\paragraph{Method.} Per-region loss weights inversely proportional to density,
$w_r = \mathrm{clip}(\bar{d}/d_r,\,1,\,4)$.
\paragraph{Result.} Gap $41 \to \mathbf{52}$. \verdictbad{Worse.}
\paragraph{Why it failed.} Increasing $w_r$ shifts the optimal operating point
downward in $\tau$; recall rises, precision falls faster, $F_1$ decreases. It
adds no information --- it is the mechanism of Eq.~\eqref{eq:f1ptau} again.

\subsection{Group-DRO}
\paragraph{Method.} Minimise the worst-group loss via soft weights
$w_g \propto \exp(\tau_{\mathrm{DRO}}\,\ell_g)$, tested at
$\tau_{\mathrm{DRO}} \in \{3, 6, 10\}$.
\paragraph{Result.} No change at any setting; gap remained $\approx 41$.
\verdictbad{Failed.}
\paragraph{Why it failed.} The Tail group's attainable $F_1$ is bounded near
$\A_{F_1}(p_{\text{tail}})$ by Eq.~\eqref{eq:af1}. Group-DRO helps when a group
is underserved because the optimiser ignores it; here the group scores low
because its events are \emph{rare}. Reallocating optimiser attention cannot
alter a base rate.

\subsection{Adaptive thresholds (first attempt)}
\paragraph{Method.} Per-region, per-category $\tau$ tuned on validation.
\paragraph{Result.} Improved to $\approx 41$, then plateaued. \verdictbad{Failed.}
\paragraph{Why it failed.} Moving a threshold slides along the precision--recall
curve without lifting it. No information enters the model.
\medskip

\noindent\fbox{\parbox{0.96\linewidth}{%
This failure proved the most important of the three. It is the first evidence
of what Section~\ref{sec:sweep} later established directly: the decision
threshold can move the fairness score a long way with no change in model
quality. Section~\ref{sec:learned} returns to it with a learned threshold.}}

%======================================================================
\section{Phase 4 --- Graph architectures}\label{sec:graphs}
\datatag{Real. Los Angeles and Chicago. 3 architectures $\times$ 4 depths $\times$ 2 cities $\times$ 5 seeds, capacity-matched, bit-reproducible.}{src/fedcrime/robust\_fair\_gnn.py --gnn \{plain,gated,attn\}}

\subsection{Method}
Let $H \in \mathbb{R}^{B\times R\times d}$ be the TCN output and $\hat{A}$ the
symmetrically normalised adjacency of a $k$-nearest-neighbour graph ($k = 6$)
built from crime-pattern correlation on the training period.
\begin{align}
\text{plain: } && H' &= \sigma(\hat{A} H W) \\
\text{gated: } && G &= \mathrm{sigmoid}(H W_g), \quad H' = \sigma\!\big(G \odot (\hat{A}HW) + (1-G)\odot H\big) \\
\text{attention: } && \alpha_{ij} &= \mathrm{softmax}_j\!\left(\tfrac{(HW)_i \cdot (HW)_j}{\sqrt{d}}\right)\Big|_{\hat A_{ij}>0}, \quad H' = \sigma(\alpha H W)
\end{align}
Depth $L \in \{1,2,3,4\}$. \textbf{Each depth is paired with its own no-graph
control at the same layer count}, so any difference is attributable to the graph
and not to capacity.

\subsection{Results --- all three variants, both cities}
Values are the graph model minus its capacity-matched control, Tail $F_1$.
Negative means the graph \emph{hurt}.
\begin{center}
\begin{tabular}{llrrrr}
\toprule
City & Variant & $L{=}1$ & $L{=}2$ & $L{=}3$ & $L{=}4$ \\
\midrule
\multirow{3}{*}{Los Angeles}
 & Plain     & $-12.92$ & $-6.51$ & $-44.49$ & $-52.82$ \\
 & Gated     & $-1.03$  & $-1.04$ & $-0.51$  & $-2.11$  \\
 & Attention & $-0.47$  & $-1.10$ & $-0.71$  & $-0.85$  \\
\midrule
\multirow{3}{*}{Chicago}
 & Plain     & $-16.29$ & $-3.55$ & $-57.52$ & $-56.01$ \\
 & Gated     & $-2.30$  & $-2.83$ & $-2.79$  & $-4.46$  \\
 & Attention & $+0.37$  & $+0.07$ & $-1.15$  & $-7.35$  \\
\bottomrule
\end{tabular}
\end{center}

\begin{center}
\begin{tabular}{lcc}
\toprule
& Los Angeles & Chicago \\
\midrule
Negative deltas & $12/12$ & $10/12$ \\
Significant at $p<0.05$ & $9/12$ & --- \\
\textbf{Significantly positive} & $\mathbf{0}$ & $\mathbf{0}$ \\
\bottomrule
\end{tabular}
\end{center}
\noindent Overall: $\mathbf{22}$ of $24$ comparisons negative, $9$ significantly
so, $\mathbf{0}$ significantly positive. The two positive Chicago values
($+0.37$, $+0.07$) are an order of magnitude below the seed-to-seed standard
deviation and are not interpreted.

\subsection{Results --- why plain GCN collapses (over-smoothing)}
Repeated application of $\hat{A}$ contracts the row space toward the leading
eigenvector. Measuring region diversity as mean pairwise distance between
$\ell_2$-normalised representations:
\begin{center}
\begin{tabular}{lc}
\toprule
Representation & Diversity \\
\midrule
Input $H$ & $1.3495$ \\
Plain GCN, $L=4$ & $\mathbf{0.0359}$ \\
\bottomrule
\end{tabular}
\end{center}
\noindent $\mathbf{97\%}$ of the distinction between regions is destroyed, after
which the model cannot assign them different predictions.

\subsection{Results --- why gated and attention do not collapse}
The gated variant is flat across depth and sits just below the control.
Inspecting the learned gate gives $G \approx 0$: offered the freedom to choose
how much neighbour information to admit, the model chooses approximately none,
and pays a small penalty for the unused parameters. Attention degrades slightly
more at depth because it can down-weight neighbours but cannot suppress them
entirely.

\medskip
\noindent\fbox{\parbox{0.96\linewidth}{%
This is the more informative of the two failures. It indicates the limitation is
\emph{informational} rather than architectural: there is little transferable
signal in the neighbourhood, and a sufficiently flexible model discovers this
unaided.}}

\subsection{Results --- Neural ODE / continuous depth}
\paragraph{Method.} Forward-Euler integration of
$\mathrm{d}h/\mathrm{d}t = \sigma(\hat{A}hW) - h$ with a \emph{single shared}
weight matrix across steps, step size $\alpha$ swept over
$\{0.05, 0.1, 0.25, 0.5, 0.75, 1.0\}$.
\paragraph{Result.} Flat at the control across the entire sweep.
\verdictbad{No improvement.} At $\alpha = 1.0$ the update reduces exactly to the
discrete GCN, confirming the parameterisation spans the intended space and that
the null result is not an implementation error.

%======================================================================
\section{Phase 5 --- The threshold sweep}\label{sec:sweep}
\datatag{Real. Chicago. No graph; only the decision threshold varies.}{analysis/threshold\_sweep.py}

\subsection{Method}
Take the model with \emph{no graph}, retrain nothing, and vary only the decision
threshold $\tau$ across $[0.05, 0.95]$ in nineteen steps.

\subsection{Results}
\begin{center}
\begin{tabular}{lccc}
\toprule
& Tail $F_1$ & says ``crime'' & Tail AUC \\
\midrule
No graph, $\tau=0.50$ (default) & $27.96$ & $22.7\%$ & $61.99$ \\
Graph model, $\tau=0.50$ & $40.83$ & $89.7\%$ & $58.94$ \\
\textbf{No graph, $\tau=0.35$} & $\mathbf{41.59}$ & $63.9\%$ & $\mathbf{61.99}$ \\
\bottomrule
\end{tabular}
\end{center}

\noindent Moving one threshold yields $\mathbf{+13.63}$ --- \emph{exceeding} the
entire architectural intervention's $+12.87$ --- with \emph{higher} AUC
($61.99$ vs $58.94$) and substantially less over-prediction ($63.9\%$ vs
$89.7\%$). AUC reads $61.99$ on \textbf{all nineteen rows} of the sweep, to the
reported precision, exactly as Eq.~\eqref{eq:auc} requires.

\medskip
\noindent\fbox{\parbox{0.96\linewidth}{%
This is the strongest single result in the project. By the metric the literature
reports, a scalar with no information content outperforms the architecture.}}

%======================================================================
\section{Phase 6 --- The physics-informed model}\label{sec:pinn}
\datatag{Real. Chicago burglary 2011--2015. $24\times24$ grid, weekly bins, 94{,}992 events, 183 training weeks, 41.5\% empty cells.}{src/pinn/build\_field.py, src/pinn/crime\_pinn.py}

\subsection{The governing equations}
We adopt the continuum limit of the Short et al.\ (2008) agent-based residential
burglary model \cite{short2008}. Two coupled fields over the city domain
$\Omega \subset \mathbb{R}^2$:
\begin{itemize}\itemsep2pt
\item $A(x,y,t)$ --- \emph{attractiveness}: the rate at which an offender
present at a location commits a burglary there;
\item $\rho(x,y,t)$ --- \emph{offender density}.
\end{itemize}

\begin{empheq}[box=\fbox]{align}
\frac{\partial A}{\partial t} &= \eta\,\nabla^2 A \;-\; A \;+\; A_0(x,y) \;+\; \rho A
\label{eq:pdeA}\\[4pt]
\frac{\partial \rho}{\partial t} &= \nabla\!\cdot\!\Big(\nabla\rho - 2\frac{\rho}{A}\nabla A\Big) \;-\; \rho A \;+\; A \;-\; A_0(x,y)
\label{eq:pderho}
\end{empheq}

\subsection{Term-by-term interpretation}
\begin{center}
\begin{tabular}{ll}
\toprule
\textbf{Term} & \textbf{Meaning} \\
\midrule
$\eta\nabla^2 A$ & attractiveness diffuses to nearby locations \\
$-A$ & it decays toward baseline over time \\
$+A_0(x,y)$ & static intrinsic appeal of a location \\
$+\rho A$ & \textbf{a crime raises local attractiveness} (repeat victimisation) \\
\midrule
$\nabla^2\rho$ & offenders move diffusively \\
$-2\nabla\!\cdot\!\big((\rho/A)\nabla A\big)$ & offenders advect \emph{up} the attractiveness gradient \\
$-\rho A$ & an offender is removed after committing a crime \\
$+A - A_0$ & replacement offenders enter the system \\
\bottomrule
\end{tabular}
\end{center}

\noindent The \textbf{observed} quantity is the crime rate
\begin{equation}
\lambda(x,y,t) \;=\; \rho(x,y,t)\,A(x,y,t).
\label{eq:obs}
\end{equation}
Neither $A$ nor $\rho$ is observed separately --- the network must infer both
from their product plus the physics constraint.

\subsection{Consistency of the transcription}
\begin{proposition}[Uniform steady state]
At a spatially uniform steady state ($\nabla A = \nabla\rho = 0$,
$\partial_t A = \partial_t \rho = 0$) both residuals vanish simultaneously
if and only if $\rho A = A - A_0$.
\end{proposition}
\begin{proof}
Eq.~\eqref{eq:pdeA} gives $0 = -A + A_0 + \rho A$, i.e.\ $\rho A = A - A_0$.
Eq.~\eqref{eq:pderho} gives $0 = -\rho A + A - A_0$, the same relation. The two
equations are therefore mutually consistent and not over-determined.
\end{proof}
\noindent \textbf{Result:} verified numerically to machine precision
($|r_A| = |r_\rho| = 0$ at $A = 1.4$, $A_0 = 0.7$, $\rho = 0.5$).

\subsection{The divergence term}
With $w \equiv \rho/A$,
\begin{equation}
\nabla\!\cdot\!\big(w\,\nabla A\big) \;=\; \nabla w \cdot \nabla A \;+\; w\,\nabla^2 A .
\label{eq:divexp}
\end{equation}
\noindent \textbf{Result:} verified against finite differences on a
$400\times400$ grid with smooth test fields; maximum relative error
$\mathbf{3.9\times10^{-5}}$.

\subsection{Loss function}
\begin{equation}
\mathcal{L} \;=\; \mathcal{L}_{\mathrm{data}} \;+\; \lambda_p\,\mathcal{L}_{\mathrm{phys}}
\end{equation}

\paragraph{Data term (Poisson).} Burglary counts $y_i \in \mathbb{Z}_{\geq0}$
are Poisson, not Gaussian. With predicted rate $\mu_i = s\,\rho_i A_i$,
\begin{equation}
\mathcal{L}_{\mathrm{data}} \;=\; \frac{1}{N}\sum_{i=1}^{N}\Big(\mu_i - y_i\log\mu_i\Big),
\label{eq:poisson}
\end{equation}
the negative Poisson log-likelihood up to a constant.

\medskip
\noindent\emph{Why this matters for fairness.} Squared error weights every cell
by $(\mu_i - y_i)^2$, so high-count (Head) cells dominate the gradient --- a
Head bias baked into the objective. \textbf{Result:} the Poisson term reduced
Tail over-prediction from $3.2\times$ to $\mathbf{1.8\times}$ the true rate.

\paragraph{Physics term.} Residuals are formed by moving all terms to one side:
\begin{align}
r_A &= \partial_t A - \big(\eta\nabla^2 A - A + A_0 + \rho A\big) \\
r_\rho &= \partial_t \rho - \big(\nabla^2\rho - 2\nabla\!\cdot\!((\rho/A)\nabla A) - \rho A + A - A_0\big)
\end{align}
\begin{equation}
\mathcal{L}_{\mathrm{phys}} \;=\; \frac{1}{N_c}\sum_{j} r_A(z_j)^2 \;+\; \frac{1}{N_c}\sum_{j} r_\rho(z_j)^2 ,
\end{equation}
evaluated at collocation points sampled \emph{inside the city mask} ---
sampling the bounding box would enforce burglary physics over Lake Michigan.
All derivatives come from automatic differentiation.

\subsection{Parameterisation}
The network $f_\theta:(x,y,t)\mapsto(A,\rho)$ uses random Fourier features
\cite{tancik2020} followed by a $5\times128$ tanh MLP:
\[
\gamma(z) = [\sin(2\pi B z),\, \cos(2\pi B z)], \qquad B \sim \mathcal{N}(0,\sigma^2),\ \sigma = 3 .
\]
Outputs pass through softplus to enforce $A > 0$ (required, since $A$ appears in
a denominator) and $\rho \geq 0$. Two physical quantities are \textbf{learned,
not fixed}: the diffusion coefficient $\eta = \mathrm{softplus}(\theta_\eta)$,
and $A_0(x,y)$ as a \emph{spatial field}, faithful to the source, which states
$A_0$ ``is not necessarily uniform over the lattice grids''
\cite{short2008,wang2020}.

\subsection{Results --- main comparison (5 seeds)}
\begin{center}
\begin{tabular}{lccccc}
\toprule
Configuration & $n$ & AUC & $F_1$ & RMSE & learned $\eta$ \\
\midrule
Short PDE + Poisson & 5 & $\mathbf{71.49\pm2.93}$ & $70.52\pm1.49$ & $1.354$ & $0.0121$ \\
Short PDE + MSE & 5 & $\mathbf{71.12\pm1.50}$ & $69.99\pm1.12$ & $1.379$ & $0.0172$ \\
No physics (control) & 5 & $53.17\pm2.46$ & $66.01\pm0.93$ & $1.750$ & --- \\
\textbf{Generic smoothness} & 1 & $\mathbf{44.63}$ & $66.90$ & $1.523$ & --- \\
\bottomrule
\end{tabular}
\end{center}

\begin{center}
\begin{tabular}{lcccc}
\toprule
vs.\ control & $\Delta$AUC & $t$ & $\Delta F_1$ & $t$ \\
\midrule
Short + Poisson & $\mathbf{+18.32}$ & $9.57^{**}$ & $+4.51$ & $5.15^{**}$ \\
Short + MSE & $+17.95$ & $12.46^{**}$ & $+3.98$ & $5.49^{**}$ \\
\bottomrule
\end{tabular}
\end{center}

\subsection{Results --- the decisive ablation}
To rule out ``any smoothness prior would do'', the PDEs are replaced by a
generic penalty containing no criminological content:
\begin{equation}
\mathcal{L}_{\mathrm{smooth}} \;=\; \frac{1}{N_c}\sum_j \big(\partial_t u\big)^2 + \big(\nabla^2 u\big)^2, \qquad u = \rho A .
\end{equation}
\noindent \textbf{Result: $\mathbf{44.63}$} --- below the no-physics control
($53.17$) \emph{and} below chance ($50$). The gain is therefore attributable to
the burglary equations specifically, not to regularisation. This is the most
reviewer-proof result in the project.

\subsection{Results --- learned diffusion constant}
$\eta \approx 0.012$--$0.017$, consistently across seeds. A small value implies
attractiveness spreads only a short distance, consistent with criminology:
near-repeat effects are measured over metres and weeks, not kilometres and
years. The model recovered a physically plausible constant it was never given.

\subsection{Results --- Tail regions (true rate 28.5\%)}
\begin{center}
\begin{tabular}{lcccc}
\toprule
Configuration & Tail $F_1$ & Tail AUC & says ``crime'' & $\times$ true rate \\
\midrule
Short + Poisson & $43.64\pm2.75$ & $\mathbf{58.35\pm3.29}$ & $52.2\%$ & $\mathbf{1.8\times}$ \\
Short + MSE & $43.36\pm3.14$ & $57.81\pm3.05$ & $56.1\%$ & $2.0\times$ \\
No physics & $43.45\pm0.98$ & $51.14\pm1.72$ & $90.2\%$ & $3.2\times$ \\
\bottomrule
\end{tabular}
\end{center}

\subsection{Results --- the measurement finding, restated}
\begin{center}
\begin{tabular}{lccc}
\toprule
Configuration & pooled AUC & Head$-$Tail $F_1$ gap & AUC gap \\
\midrule
Short + Poisson & $71.49$ & $47.32\pm2.65$ & $\mathbf{-1.92}$ \\
Short + MSE & $71.12$ & $47.60\pm3.21$ & $\mathbf{-3.50}$ \\
No physics & $53.17$ & $45.54\pm1.91$ & $+0.33$ \\
\midrule
\textbf{Range} & $\mathbf{18.32}$ & $\mathbf{2.06}$ & --- \\
\bottomrule
\end{tabular}
\end{center}
\noindent Real skill spans $18.32$ points; the reported fairness gap spans
$2.06$. Sharper still, \emph{within the Tail group alone}, $F_1$ reads
$43.64 / 43.36 / 43.45$ --- a spread of $\mathbf{0.28}$ --- while Tail AUC spans
$\mathbf{7.2}$. Same group, same base rate, same threshold procedure.

\medskip
\noindent \textbf{The sign also reverses.} The $F_1$ gap reports Head regions as
$\sim$47 points better served; the AUC gap reports Head as \emph{worse} served
in both physics models. The two metrics disagree about \emph{who is being
failed}.

%======================================================================
\section{Phase 7 --- Learned thresholds}\label{sec:learned}
\datatag{\textbf{Synthetic.} Three groups; $P(s\mid y)$ identical in all groups so skill is equal by construction; base rates $0.60 / 0.35 / 0.12$. See Section~\ref{sec:phases1} for why real data cannot answer this question.}{learned\_thresholds.py}

\subsection{Motivation}
The supervisor proposed \emph{adaptive} thresholds: not a fixed rule, and not a
hand-written formula, but a threshold the model determines from data. This
returns to the intervention of Section~\ref{sec:interventions} with the
threshold made learnable.

\subsection{Method}
A threshold is a step function and has no gradient. We replace the hard
comparison with a sigmoid relaxation, which is differentiable:
\begin{equation}
p_i = \mathrm{sigmoid}\big(\beta (s_i - \tau)\big), \qquad \beta = 50,
\end{equation}
\begin{align}
TP = \textstyle\sum_i p_i y_i, \quad
FP = \textstyle\sum_i p_i (1-y_i), \quad
FN = \textstyle\sum_i (1-p_i) y_i,
\end{align}
\begin{equation}
\text{soft-}F_1 = \frac{2\,TP}{2\,TP + FP + FN}.
\end{equation}
At $\beta = 50$ the sigmoid closely approximates a step, so the soft counts
track the true counts while retaining a gradient. Two objectives were compared:
\begin{itemize}\itemsep2pt
\item \textbf{Arm A} --- maximise soft-$F_1$ within each group;
\item \textbf{Arm B} --- match each group's false-positive rate $F_g$ to the
pooled rate $\varphi$ at $\tau = 0.50$: $\;\mathcal{L} = \sum_g (F_g - \varphi)^2$.
\end{itemize}

\subsection{Results}
\begin{center}
\begin{tabular}{lccc}
\toprule
Method & $F_1$ gap & AUC gap & \textbf{TPR gap} \\
\midrule
Fixed $\tau = 0.50$ & $+37.66$ & $-0.07$ & $-0.13$ \\
\textbf{Arm A} --- maximise soft-$F_1$ & $+35.14$ & $-0.07$ & $\mathbf{+35.56}$ \\
\textbf{Arm B} --- matched FPR & $+37.15$ & $-0.07$ & $-1.57$ \\
\bottomrule
\end{tabular}
\end{center}

\subsection{Results --- what Arm A did, per group}
\begin{center}
\begin{tabular}{lcc}
\toprule
Group & Threshold chosen & TPR (\% of real events caught) \\
\midrule
Head & $0.425$ & $89.10$ \\
Mid & $0.520$ & $75.34$ \\
Tail & $\mathbf{0.601}$ & $\mathbf{53.54}$ \\
\bottomrule
\end{tabular}
\end{center}

\noindent Arm A set the \emph{highest} threshold for the sparsest group,
producing a $35.56$-point disparity in who is caught --- on data where every
group is equally predictable by construction.
\verdictbad{Arm A is actively harmful.} It manufactures precisely the harm the
intervention was intended to remove.

\subsection{Why Arm A chose a higher threshold for the sparse group}
The classifier is calibrated to the \emph{pooled} base rate, not the per-group
rate. In a low-base-rate group its scores systematically overstate
$\Pr(y=1)$, so the threshold must rise to compensate. This is arithmetic, not
prejudice --- but the distributional consequence is identical.

\medskip
\noindent \emph{Relation to prior work.} Lipton et al.\ (2014)
\cite{lipton2014} show the $F_1$-optimal threshold equals $F_1^*/2$ for
\emph{calibrated} probabilities, which predicts a \emph{lower} threshold for a
sparse group. Our scores are calibrated to the pooled rate, not per group, which
accounts for the opposite direction. Any write-up must engage with this.

\subsection{Interpretation}
\verdictgood{Arm B is correct but adds no skill.} AUC moved by $0.07$ in all
three conditions --- the model was never altered. The conclusion is that
adaptive thresholding is \textbf{not a remedy}; it is a \textbf{measuring
instrument} that quantifies how far the reported score can travel while the
model stands still.

\subsection{An error we made and corrected}
Our first Arm B objective minimised the spread in \emph{true}-positive rate.
That objective is degenerate: it drove every threshold to $0.02$, so the model
predicted positive everywhere and every group attained $\mathrm{TPR} = 100\%$
--- a perfect fairness score from a useless model. This was caught in an
independent NumPy validation before any result was trusted, and the objective
was rewritten to match false-positive rates.

%======================================================================
\section{Phase 8 --- Auditing the fairness toolkits}\label{sec:tools}
\datatag{\textbf{Synthetic.} $s \mid y{=}1 \sim \mathcal{N}(1,1)$ and $s \mid y{=}0 \sim \mathcal{N}(0,1)$, identical in all three groups; base rates $0.60/0.35/0.12$. Fairlearn 0.14.0, AIF360 0.6.1, scikit-learn 1.7.2.}{audit\_fairness\_tools.py}

\subsection{Motivation}
The preceding sections concern our own models. The question of external
relevance is whether the same failure appears in the tooling practitioners
actually deploy. Fairlearn (Microsoft) and AIF360 (IBM) are the two most widely
used fairness libraries; both ship post-processing methods that adjust per-group
decision thresholds.

\subsection{Validation of the test bed}
Before invoking any library, equal skill was verified by hand: the TPR of each
group at a matched false-positive rate.
\begin{center}
\begin{tabular}{lccc}
\toprule
FPR held at & Head & Mid & Tail \\
\midrule
$5\%$ & $30.2$ & $25.9$ & $29.9$ \\
$10\%$ & $41.8$ & $38.4$ & $40.5$ \\
$20\%$ & $56.7$ & $57.1$ & $56.5$ \\
\bottomrule
\end{tabular}
\end{center}
\noindent Equal within sampling noise, as constructed. On this data the correct
audit finding is that no group is served worse.

\subsection{Results --- the metric functions}
\begin{center}
\begin{tabular}{lccccc}
\toprule
Group & accuracy & precision & recall & $F_1$ & AUC \\
\midrule
Head & $63.46$ & $84.84$ & $48.00$ & $61.31$ & $76.87$ \\
Mid & $72.99$ & $64.12$ & $47.07$ & $54.28$ & $76.24$ \\
Tail & $81.53$ & $32.16$ & $45.82$ & $37.79$ & $76.03$ \\
\midrule
\textbf{\texttt{MetricFrame.difference()}} & $18.06$ & $\mathbf{52.67}$ & $2.18$ & $\mathbf{23.52}$ & $\mathbf{0.84}$ \\
\bottomrule
\end{tabular}
\end{center}
\noindent Fairlearn reports a $52.67$-point precision disparity and a
$23.52$-point $F_1$ disparity on groups with identical ROC curves. Recall
($2.18$) and AUC ($0.84$) are correct.

\subsection{Results --- Fairlearn \texttt{ThresholdOptimizer}, 5 seeds}
\begin{center}
\begin{tabular}{lccc}
\toprule
\texttt{constraints} & flagged gap & \textbf{TPR gap} & AUC gap \\
\midrule
baseline ($\tau = 0.50$) & $+15.98\pm0.66$ & $+0.35\pm1.32$ & $+0.16\pm0.68$ \\
\texttt{demographic\_parity} \textbf{[DEFAULT]} & $-0.02\pm0.16$ & $\mathbf{-21.26\pm0.52}$ & $+0.16\pm0.68$ \\
\texttt{true\_positive\_rate\_parity} & $+15.93\pm0.74$ & $-0.02\pm0.12$ & $+0.16\pm0.68$ \\
\texttt{false\_positive\_rate\_parity} & $+15.66\pm1.16$ & $+0.27\pm2.23$ & $+0.16\pm0.68$ \\
\texttt{equalized\_odds} & $+15.21\pm0.78$ & $+0.18\pm0.32$ & $+0.16\pm0.68$ \\
\bottomrule
\end{tabular}
\end{center}

\noindent The library default \emph{opens} a $21.26$-point TPR disparity where
the pre-intervention gap was $+0.35 \pm 1.32$. Per group under the default:
\begin{center}
\begin{tabular}{lccc}
\toprule
Group & flagged & TPR & FPR \\
\midrule
Head & $24.2\%$ & $35.50$ & $7.12$ \\
Mid & $24.3\%$ & $46.07$ & $13.05$ \\
Tail & $24.3\%$ & $56.33$ & $\mathbf{19.84}$ \\
\bottomrule
\end{tabular}
\end{center}
\noindent Selection rates are equalised as promised, but the sparse group's
false-alarm rate is now $2.8\times$ the dense group's. In a policing deployment
that is wrongful surveillance of the low-crime area, produced by the fairness
intervention.

\subsection{Results --- AIF360 post-processors}
Head privileged, Tail unprivileged; AIF360 reports unprivileged minus
privileged.
\begin{center}
\begin{tabular}{lcccc}
\toprule
Method & DP diff & EO diff & Tail TPR & \textbf{Tail FPR} \\
\midrule
baseline ($\tau = 0.50$) & $-16.68$ & $-2.18$ & $45.82$ & $13.49$ \\
\texttt{EqOddsPostprocessing} & $-15.54$ & $+0.18$ & $45.82$ & $13.49$ \\
\texttt{CalibratedEqOdds} (fnr) & $-16.68$ & $-2.18$ & $45.82$ & $13.49$ \\
\texttt{RejectOption} (stat.\ parity) & $-4.46$ & $+12.26$ & $78.06$ & $\mathbf{41.55}$ \\
\texttt{RejectOption} (equal opp.) & $-20.12$ & $-2.75$ & $68.67$ & $31.54$ \\
\bottomrule
\end{tabular}
\end{center}
\noindent \texttt{RejectOptionClassification} on statistical parity raises the
Tail group's false-positive rate from $13.49\%$ to $\mathbf{41.55\%}$ --- three
times as many crime-free locations flagged --- to close a parity gap that was an
artifact of the base rate. \texttt{CalibratedEqOdds} changed nothing at all.
Per-group AUC is $76.9 / 76.0$ on every row: post-processing re-thresholds, it
does not re-rank.

\subsection{Scope of this claim --- stated precisely}
The constraints targeting TPR or both error rates behave \emph{correctly} here,
with gaps near zero. The failure is specific to the parity-of-selection family,
which happens to contain the default.

\medskip
\noindent\fbox{\parbox{0.96\linewidth}{%
\textbf{Honesty requirement.} The incompatibility of demographic parity and
equalized odds under unequal base rates is a \emph{known} theoretical result
\cite{kleinberg2016,chouldechova2017} and must be cited as such. Our
contribution is narrower and must be stated as: this known incompatibility is
the \emph{default configuration} of widely deployed software, and its cost is
quantifiable ($21.26$-point TPR gap; $3\times$ wrongful flagging).}}

%======================================================================
\section{Experimental protocol}\label{sec:protocol}

\subsection{Determinism}
Seeding \texttt{torch} alone is insufficient on GPU: cuDNN selects algorithms by
autotuning and several CUDA kernels accumulate non-deterministically. Our
initial results were consequently irreproducible (Section~\ref{sec:mistakes}).
All reported results use:
\begin{lstlisting}
def seed_everything(seed: int):
    os.environ.setdefault("CUBLAS_WORKSPACE_CONFIG", ":4096:8")
    random.seed(seed); np.random.seed(seed)
    torch.manual_seed(seed); torch.cuda.manual_seed_all(seed)
    torch.backends.cudnn.deterministic = True
    torch.backends.cudnn.benchmark     = False
    torch.use_deterministic_algorithms(True, warn_only=True)
\end{lstlisting}
\noindent plus gradient clipping at norm $1.0$ and a non-finite-loss guard. A
\texttt{-{}-repro-check} mode trains an identical configuration twice and
asserts bit-identical output. \textbf{Result: difference $0.00000000$.}

\subsection{Capacity-matched controls}
\begin{lstlisting}
for depth in [1, 2, 3, 4]:
    m_graph   = run(use_graph=True,  n_layers=depth)   # treatment
    m_control = run(use_graph=False, n_layers=depth)   # same size, no graph
    delta = m_graph["overall"] - m_control["overall"]  # isolates the graph
\end{lstlisting}

\subsection{Statistics}
Five seeds per configuration; mean $\pm$ standard deviation reported. Group
comparisons use the Welch $t$-statistic
\begin{equation}
t \;=\; \frac{\bar{x}_1 - \bar{x}_2}{\sqrt{s_1^2/n_1 + s_2^2/n_2}},
\end{equation}
with $|t| > 2.78$ ($df = 4$) taken as $p < 0.05$. The literature audit
additionally uses permutation tests ($2\times10^5$ permutations), Fisher's
method for combining $p$-values, and a bootstrap CI over models as the unit of
generalisation.

%======================================================================
\section{Summary of every method tested}

\begin{center}
\begin{tabular}{clccl}
\toprule
\# & Method & Data & Effect vs.\ matched control & Verdict \\
\midrule
1 & Reweighting & real & gap $41 \to 52$ & \verdictbad{worse} \\
2 & Group-DRO ($\tau = 3,6,10$) & real & no change & \verdictbad{fail} \\
3 & Adaptive thresholds & real & plateau at $\sim41$ & \verdictbad{fail} \\
4 & Plain GCN, $L = 1\ldots4$ & real & $-12.9$ to $-57.5$ & \verdictbad{worse} \\
5 & Gated GCN, $L = 1\ldots4$ & real & $-0.5$ to $-4.5$ & \verdictbad{worse} \\
6 & Attention GCN, $L = 1\ldots4$ & real & $-0.5$ to $-7.4$ & \verdictbad{worse} \\
7 & Neural ODE, $\alpha$ swept & real & flat at control & \verdictbad{fail} \\
8 & \textbf{PINN (Short et al.\ PDE)} & real & $\mathbf{+18.32}$ AUC, $t = 9.57$ & \verdictgood{works} \\
9 & Learned thresholds, Arm A & synth. & $+35.56$ TPR gap created & \verdictbad{harmful} \\
10 & Learned thresholds, Arm B & synth. & correct, no skill added & \verdictgood{diagnostic} \\
\bottomrule
\end{tabular}
\end{center}

\noindent One intervention of ten added genuine skill. Five demonstrations ---
the graph experiments, the threshold sweep, the PINN configurations, the learned
thresholds, and the library audit --- independently establish the same
measurement finding.

%======================================================================
\section{Errors identified and corrected}\label{sec:mistakes}
\begin{center}
\begin{tabular}{p{5cm}p{9.8cm}}
\toprule
\textbf{Error} & \textbf{Resolution} \\
\midrule
Results not reproducible; identical commands gave $48.86$ and $54.03$ & Only \texttt{torch.manual\_seed} was set. Fixed by full determinism (Section~\ref{sec:protocol}). \textbf{All prior model results discarded.} \\
\addlinespace
Controls not capacity-matched: depth-4 graph vs.\ 1-layer control & Every depth now paired with its own control at the same layer count. \\
\addlinespace
An apparent fairness breakthrough ($+16.95$ Tail $F_1$, corrected gap $\to -2.39$) & Tail AUC had \emph{fallen} $4.01$ and positive-prediction rate rose $23\% \to 80\%$. Withdrawn as a positive result; became the central finding. \\
\addlinespace
All per-group metrics returned NaN & Labels were derived from the smoothed field, making every label 1. Fixed to use raw counts. \\
\addlinespace
Collocation sampled from the bounding box & Physics was enforced over Lake Michigan. Fixed to sample within the city mask. \\
\addlinespace
Degenerate Arm B objective drove all thresholds to $0.02$ & Minimising TPR spread rewards predicting positive everywhere. Rewritten as matched FPR (Section~\ref{sec:learned}). \\
\addlinespace
\texttt{-{}-sampling adaptive} was described as Lin \& Chen's method & Theirs is causality-guided; ours was plain residual-based. Documented as not implemented. \\
\bottomrule
\end{tabular}
\end{center}
\noindent We also corrected claims in our own writing: $\A_{F_1}$ is the
\emph{skill-free baseline}, not a ceiling (models exceed it); the theory of
Section~\ref{sec:theory} is established, not novel; and bit-reproducibility
holds within, not across, hardware.

\paragraph{Excluded work.} A sparsity-camouflaged attack and a density/graph
``vouching'' aggregator were developed in an earlier phase. Those results
predate the determinism fix and are \textbf{void}; they are retained in the
repository for continuity and excluded from any write-up.

%======================================================================
\section{Limitations}
\begin{enumerate}\itemsep3pt
\item \textbf{The PINN learns \emph{where}, not \emph{when}.} Per-group AUC is
$51$--$58$; most of the pooled $71.49$ derives from discriminating between
neighbourhoods rather than between weeks within one.
\item \textbf{Residual over-prediction.} Even the best configuration flags
$1.8\times$ the true Tail rate.
\item \textbf{Scale.} Two cities, one country; one crime type for the PINN.
\item \textbf{The region graph is correlational}, built from crime-pattern
similarity rather than physical street adjacency. A street-network graph is a
natural further test.
\item \textbf{Sections~\ref{sec:learned} and \ref{sec:tools} use synthetic
data.} This is required to establish their claims, but it means we show these
methods \emph{can} be fooled, not how severely they fail on Chicago
specifically. Running the learned-threshold experiment on real Chicago scores
remains outstanding.
\item \textbf{Recorded crime is not crime.} It measures police activity. Base
rates are themselves shaped by policing, so the confound compounds.
\item \textbf{The burglary model derives from Broken Windows theory}, which is
contested \cite{sampson2004,goodson2021}. We implement it in order to
\emph{measure} its behaviour, particularly in low-crime neighbourhoods; we do
not endorse it as a description of crime.
\item Cross-hardware variation affects the third significant digit.
\end{enumerate}

%======================================================================
\section{Recommendations}
\begin{enumerate}\itemsep2pt
\item Never report a threshold-dependent fairness gap alone.
\item Base-rate correction is \emph{not} sufficient; state the operating point.
\item Report an invariant metric (AUC, balanced accuracy) alongside any $F_1$ gap.
\item Report the positive-prediction rate per group.
\item Treat any gap closure unmatched by an invariant metric as \textbf{unverified}.
\item When using a fairness library, do not accept the default constraint
without checking what it optimises.
\end{enumerate}

%======================================================================
\begin{thebibliography}{99}\small
\bibitem{boyd2012} K. Boyd, V. S. Costa, J. Davis, C. D. Page.
\emph{Unachievable region in precision-recall space and its effect on empirical evaluation}. ICML, 2012.
\bibitem{chouldechova2017} A. Chouldechova.
\emph{Fair prediction with disparate impact: a study of bias in recidivism prediction instruments}. Big Data 5(2), 2017.
\bibitem{davis2006} J. Davis, M. Goadrich.
\emph{The relationship between precision-recall and ROC curves}. ICML, 2006.
\bibitem{ensign2018} D. Ensign, S. A. Friedler, S. Neville, C. Scheidegger, S. Venkatasubramanian.
\emph{Runaway feedback loops in predictive policing}. FAccT, 2018.
\bibitem{goodson2021} H. Goodson, A. Hoyer-Leitzel.
\emph{Examining the modeling framework of crime hotspot models in predictive policing}. arXiv:2103.11757, 2021.
\bibitem{kleinberg2016} J. Kleinberg, S. Mullainathan, M. Raghavan.
\emph{Inherent trade-offs in the fair determination of risk scores}. ITCS, 2017 (arXiv 2016).
\bibitem{lipton2014} Z. C. Lipton, C. Elkan, B. Naryanaswamy.
\emph{Optimal thresholding of classifiers to maximize F1 measure}. ECML PKDD, 2014.
\bibitem{lum2016} K. Lum, W. Isaac.
\emph{To predict and serve?} Significance 13(5), 2016.
\bibitem{richardson2019} R. Richardson, J. Schultz, K. Crawford.
\emph{Dirty data, bad predictions}. NYU Law Review, 2019.
\bibitem{sampson2004} R. J. Sampson, S. W. Raudenbush.
\emph{Seeing disorder: neighborhood stigma and the social construction of broken windows}. Social Psychology Quarterly 67(4), 2004.
\bibitem{short2008} M. B. Short, M. R. D'Orsogna, V. B. Pasour, G. E. Tita, P. J. Brantingham, A. L. Bertozzi, L. B. Chayes.
\emph{A statistical model of criminal behavior}. Math.\ Models Methods Appl.\ Sci.\ 18(supp01), 2008.
\bibitem{tancik2020} M. Tancik et al.
\emph{Fourier features let networks learn high frequency functions in low dimensional domains}. NeurIPS, 2020.
\bibitem{wang2020} C. Wang, Y. Zhang, A. L. Bertozzi, M. B. Short.
\emph{A stochastic-statistical residential burglary model with independent Poisson clocks}. European J.\ Applied Mathematics, 2020.
\end{thebibliography}

\end{document}
